% Corrected French for printed p.124.
% Correction: injections canonique -> injections canoniques.
\SourcePage{124}

\LayerTarget{tohoku:section:1.1:p8}
\noindent naturel de choisir alors parmi tous les $(A,(u_i))$ comme ci-dessus,
un système particulier (par exemple au moyen du symbole $\tau$ de Hilbert)
qu'on appellera \textit{le produit} de la famille d'objets $(A_i)_{i\in I}$.
C'est donc, non un simple objet $A$ de $\mathbf{C}$, mais un tel objet muni d'une
famille $(u_i)$ de morphismes dans les $A_i$, appelés les
\textit{projections canoniques du produit} sur ses facteurs $A_i$. On note le
produit des $A_i$ (s'il existe) par $\displaystyle\prod_{i\in I}A_i$. Si $I$ est réduit à
un élément $i$, alors le produit s'identifie à $A_i$ lui même. On dit que $\mathbf{C}$
est une \textit{catégorie avec produits} si le produit de deux objets de $\mathbf{C}$
existe toujours (alors, il en est de même du produit d'une famille finie non
vide quelconque d'objets de $\mathbf{C}$). On dit que $\mathbf{C}$ est une \textit{catégorie
avec produits infinis} si le produit d'une famille non vide quelconque d'objets
de $\mathbf{C}$ existe toujours. Nous avons vu que si on a deux produits
$A=\displaystyle\prod_{i\in I}A_i$ et $B=\displaystyle\prod_{i\in I}B_i$ correspondants à un même
ensemble d'indices $I$, alors une famille $(v_i)$ de morphismes
$A_i\to B_i$ définit canoniquement un morphisme $v$ de $A$ dans
$B$, appelé \textit{produit des morphismes $v_i$} et parfois noté
$\displaystyle\prod_{i\in I}v_i$. Si les $v_i$ sont des monomorphismes, il en est de même
de leur produit, mais l'énoncé analogue pour les épimorphismes n'est pas vrai
en général, (comme on voit par exemple sur la catégorie des faisceaux sur un
espace topologique fixé).

\LayerTarget{tohoku:section:1.1:p9}
Des considérations duales des précédentes permettent de définir la notion de
\textit{représentation d'un objet $A$ comme somme} d'une famille d'objets $A_i$
par des morphismes $u_i:A_i\to A$ (pour tout $B\in \mathbf{C}$,
l'application naturelle
\[
  \mathrm{Hom}(A,B)\to
  \prod_{i\in I}\mathrm{Hom}(A_i,B)
\]
est bijective), de \textit{somme directe} $\displaystyle\bigoplus_{i\in I}A_i$, munie des
\textit{injections canoniques} $A_i\to\displaystyle\bigoplus_{i\in I}A_i$ (qui
d'ailleurs ne sont pas nécessairement des monomorphismes, malgré leur nom), de
\textit{morphisme somme} d'une famille de morphismes
$u_i:A_i\to B_i$. Si les $u_i$ sont des épimorphismes, il en est de
même de leur somme.

\LayerTarget{tohoku:section:1.2}
\noindent\textbf{1.2. Foncteurs.}\quad
\LayerTarget{tohoku:section:1.2:p1}
Soient $\mathbf{C},\mathbf{C}'$ deux catégories. Rappelons qu'on appelle \textit{foncteur
covariant} de $\mathbf{C}$ dans $\mathbf{C}'$ une ``fonction'' $F$ qui à un objet $A\in \mathbf{C}$,
associe un objet $F(A)$ de $\mathbf{C}'$, et à un morphisme
$u:A\to B$ dans $\mathbf{C}$, un morphisme
$F(u):F(A)\to F(B)$, de telle façon qu'on ait
$F(I_A)=I_{F(A)}$ et $F(vu)=F(v)F(u)$. Définition analogue pour les
\textit{foncteurs contravariants} de $\mathbf{C}$ dans $\mathbf{C}'$ (qui sont aussi des
foncteurs covariants de $\mathbf{C}^{\circ}$ dans $\mathbf{C}'$ ou de $\mathbf{C}$ dans
${\mathbf{C}'}^{\circ}$). On définit de même les foncteurs de plusieurs variables ou
\textit{multifoncteurs}, covariants en certaines variables et contravariants en
d'autres. Dans les généralités, nous nous bornerons pour simplifier aux
foncteurs d'une variable. Les foncteurs se composent comme des fonctions,
cette composition est associative et les ``foncteurs identiques'' jouent le
rôle d'unités.

\LayerTarget{tohoku:section:1.2:p2}
Soient $\mathbf{C}$ et $\mathbf{C}'$ deux catégories fixées, $F$ et $G$ deux foncteurs covariants
de $\mathbf{C}$ dans $\mathbf{C}'$, un \textit{morphisme fonctoriel} $f$ de $F$ dans $G$ (appelé
aussi ``transformation

\TohokuContinuation{125}
